\subsection{Pontryagin duality theorem}

THEOREM 2 (Pontryagin). --- The canonical mapping η from G into \(\widehat{G}\) is an isomorphism of topological groups. It maps the Haar measure dx to the bi-dual Haar measure d \(\widehat{x}\).

Let us first prove that \(\eta\) is injective and strict. It suffices for this to show that for every neighborhood U of \(e\) in G, there exists a neighborhood W of \(e\) in \(\widehat{\mathbf{G}}\) such that \(\overline{\eta}^{-1}(\mathrm{W})\subset \mathrm{U}\) (lemma 2 of II, p.~200). Now let V be a compact symmetric neighborhood of \(e\) in G such that \(\mathrm{V}^2\subset \mathrm{U}\), let \(f\) be a continuous positive function on G, with support contained in V, and such that \(f(e) > 0\). Let \(g = \widetilde{f}*f\). Then \(g\) belongs to \(\mathrm{A}(\mathrm{G})\), its support is contained in U, and \(g(e) > 0\). Moreover, \(\mathcal{F}_{\mathrm{G}}(g)\in \mathrm{L}^{1}(\widehat{\mathrm{G}})\) by prop. 11 of II, p.~217. The set W of \(\xi\) in \(\widehat{\mathbf{G}}\) such that

\[
\left| \overline {{\mathcal {F}}} _ {\widehat {\mathrm{G}}} (\mathcal {F} _ {\mathrm{G}} (g)) (\xi) - \overline {{\mathcal {F}}} _ {\widehat {\mathrm{G}}} (\mathcal {F} _ {\mathrm{G}} (g)) (e) \right| <   \frac {1}{2} g (e)
\]

is a neighborhood of \(e\) in \(\widehat{\mathbf{G}}\) since the function \(\overline{\mathcal{F}}_{\widehat{\mathrm{G}}}(\mathcal{F}_{\mathrm{G}}(g))\) is continuous on \(\widehat{\mathbf{G}}\) . Let \(x \in \overline{\eta}^{-1}(\mathrm{W})\) . By formula (26), we have

\[
\overline {{\mathcal {F}}} _ {\widehat {\mathrm{G}}} (\mathcal {F} _ {\mathrm{G}} (g)) (\eta (x)) = g (x),
\]

and therefore \(|g(x) - g(e)| < \frac{1}{2} g(e)\) . This implies \(g(x) \neq 0\) and therefore \(x \in \mathrm{U}\) , since the support of \(g\) is contained in U. Thus \(\overline{\eta}^{-1}(\mathrm{W}) \subset \mathrm{U}\) .

Let us prove that the map \(\eta\) is surjective. Since this map is a homeomorphism onto its image, the group \(\eta(\mathbf{G})\) is a locally compact subgroup of \(\widehat{\widehat{\mathbf{G}}}\). It is therefore closed in \(\widehat{\widehat{\mathbf{G}}}\) (TG, III, p.~22, cor. 2). Let us argue by contradiction and suppose that there exists a character \(\xi \in \widehat{\widehat{\mathbf{G}}}\) such that \(\xi \notin \eta(\mathbf{G})\). There then exists (corollary 1 of II, p.~219) a non-zero element \(f\) of \(\mathrm{L}^1(\widehat{\mathrm{G}})\) such that \(\mathcal{F}_{\widehat{\mathrm{G}}} (f)\) vanishes on \(\eta(\mathrm{G})\). Let \(g \in \mathrm{L}^1(\mathrm{G})\). The function \((x, \chi) \mapsto g(x)f(\chi)\overline{\langle\chi,x\rangle}\) belongs to \(\mathrm{L}^1(\mathrm{G} \times \widehat{\mathrm{G}})\). By the Lebesgue-Fubini theorem (INT, V, §8, no. 4, th. 1, a)), it follows that

\[
\begin{array}{r} \int_ {\widehat {\mathrm{G}}} f (\chi) \mathcal {F} _ {\mathrm{G}} (g) (\chi) d \chi = \int_ {\mathrm{G}} g (x) \Bigl (\int_ {\widehat {\mathrm{G}}} f (\chi) \overline {{\langle \chi , x \rangle}} d \chi \Bigr) d x \\ = \int_ {\mathrm{G}} g (x) \mathcal {F} _ {\widehat {\mathrm{G}}} (f) (\eta (x)) d x = 0. \end{array}
\]

Since the image of the Fourier transform is dense in \(\mathcal{C}_{0}(\widehat{\mathrm{G}})\) (cor. of prop. 5 of II, p.~209), it follows that the measure \(f \cdot d\chi\) is zero. This contradicts the fact that \(f \neq 0\) in \(\mathrm{L}^{1}(\widehat{\mathrm{G}})\), and proves that \(\eta\) is surjective.

The image measure \(\eta(dx)\) and the measure \(\nu\) dual to the measure \(d\chi\) are Haar measures on \(\widehat{\widehat{G}}\) . Let \(f\) be a non-zero element of A(G); in particular \(f \in \mathrm{L}^2(\mathrm{G})\) . By prop. 12 of II, p.~219, \(\mathcal{F}_{\mathrm{G}}(f) \in \mathrm{L}^{1}(\widehat{\mathrm{G}})\) and one has

\[
\int_ {\widehat {\widehat {\mathrm{G}}}} \Big | \overline {{\mathcal {F}}} _ {\widehat {\mathrm{G}}} (\mathcal {F} _ {\mathrm{G}} (f)) \Big | ^ {2} \eta (d x) = \int_ {\mathrm{G}} | f | ^ {2} d x = \int_ {\widehat {\widehat {\mathrm{G}}}} \Big | \overline {{\mathcal {F}}} _ {\widehat {\mathrm{G}}} (\mathcal {F} _ {\mathrm{G}} (f)) \Big | ^ {2} d \nu ,
\]

where the second equality follows from two applications of the Plancherel formula, so the dual Haar measure of \(d\chi\) is the measure \(\eta(dx)\) .

We shall henceforth identify G and \(\widehat{\mathrm{G}}\) by the isomorphism \(\eta\) . We then have:

COROLLARY. --- The Fourier cotransform of \(L^{2}(\widehat{G})\) on \(L^{2}(G)\) and the Fourier transform of \(L^{2}(G)\) on \(L^{2}(\widehat{G})\) are reciprocal isometries of each other.

Remark. --- Let \(f \in \mathrm{L}^2(\mathrm{G})\) and \(g \in \mathrm{L}^2(\widehat{\mathrm{G}})\) . Applying the Plancherel formula (24) to \(f\) and \(\overline{\mathcal{F}_{\widehat{\mathrm{G}}}(g)}\) , we obtain the formula

\[
\int_ {\mathrm{G}} f (x) \mathcal {F} _ {\widehat {\mathrm{G}}} (g) (x) d x = \int_ {\widehat {\mathrm{G}}} \mathcal {F} _ {\mathrm{G}} (f) (\chi) g (\chi) d \widehat {x} (\chi)\tag{29}
\]

since we have \(\mathcal{F}_{\mathrm{G}}(\overline{\mathcal{F}_{\widehat{\mathrm{G}}}(\boldsymbol{g})}) = \mathcal{F}_{\mathrm{G}}(\overline{\mathcal{F}}_{\widehat{\mathrm{G}}}(\overline{\boldsymbol{g}})) = \overline{\boldsymbol{g}}.\)

The Fourier transform and cotransform defined on \(\mathcal{M}^{1}(\widehat{\mathrm{G}})\) take values in the space of continuous bounded functions on G. For \(\beta\in\mathcal{M}^{1}(\widehat{\mathrm{G}})\) and \(x\in\mathrm{G}\) , one has

\[
\mathcal {F} _ {\widehat {\mathrm{G}}} (\beta) (x) = \int_ {\widehat {\mathrm{G}}} \overline {{\langle \chi , x \rangle}} d \beta (\chi), \qquad \overline {{\mathcal {F}}} _ {\widehat {\mathrm{G}}} (\beta) (x) = \int_ {\widehat {\mathrm{G}}} \langle \chi , x \rangle d \beta (\chi).
\]

The Fourier transforms of G and \(\widehat{G}\) are also transposes of each other. More precisely:

PROPOSITION 13. --- Let \(\alpha \in \mathcal{M}^{1}(\mathrm{G})\) and \(\beta \in \mathcal{M}^{1}(\widehat{\mathrm{G}})\) . We then have

\[
\mathcal {F} _ {\mathrm{G}} (\overline {{\mathcal {F}}} _ {\widehat {\mathrm{G}}} (\beta) \cdot \alpha) = \beta * \mathcal {F} _ {\mathrm{G}} (\alpha)\tag{30}
\]

and in particular

\[
\int_ {\mathrm{G}} \mathcal {F} _ {\widehat {\mathrm{G}}} (\beta) (x) d \alpha (x) = \int_ {\widehat {\mathrm{G}}} \mathcal {F} _ {\mathrm{G}} (\alpha) (\chi) d \beta (\chi).\tag{31}
\]

Formula (30) implies formula (31) by evaluating both sides of the identity at \(\chi = 1\) . Let us prove (30). Let \(\chi \in \widehat{G}\). It follows that

\[
\begin{array}{r l} & {(\mathcal {F} _ {\mathrm{G}} (\overline {{\mathcal {F}}} _ {\widehat {\mathrm{G}}} (\beta) \cdot \alpha)) (\chi) = \int_ {\mathrm{G}} \overline {{\langle \chi , x \rangle}} \overline {{\mathcal {F}}} _ {\widehat {\mathrm{G}}} (\beta) (x) d \alpha (x)} \\ & {\qquad = \int_ {\mathrm{G}} \overline {{\langle \chi , x \rangle}} \Big (\int_ {\widehat {\mathrm{G}}} \langle \xi , x \rangle d \beta (\xi) \Big) d \alpha (x).} \end{array}
\]

The function \((x,\xi)\mapsto\overline{\langle\chi,x\rangle}\langle\xi,x\rangle\) is continuous and bounded, hence integrable on \(G\times\widehat{G}\) with respect to the measure \(\alpha\otimes\beta\) . By the Lebesgue-Fubini theorem (INT, V, §8, no. 4, th. 1, a)), we obtain

\[
\begin{array}{r l} & {(\mathcal {F} _ {\mathrm{G}} (\overline {{\mathcal {F}}} _ {\widehat {\mathrm{G}}} (\beta) \cdot \alpha)) (\chi) = \int_ {\widehat {\mathrm{G}}} \Bigl (\int_ {\mathrm{G}} \overline {{\langle \chi \xi^ {- 1} , x \rangle}} d \alpha (x) \Bigr) d \beta (\xi)} \\ & {\qquad = \int_ {\widehat {\mathrm{G}}} \mathcal {F} _ {\mathrm{G}} (\alpha) (\chi \xi^ {- 1}) d \beta (\xi) = (\beta * \mathcal {F} _ {\mathrm{G}} (\alpha)) (\chi),} \end{array}
\]

as desired.

COROLLARY. --- The Fourier transform \(\mathcal{F}_{\mathrm{G}}\) is injective on \(\mathcal{M}^{1}(\mathrm{G})\) .

Indeed, if \(\alpha \in \mathcal{M}^1 (\mathrm{G})\) satisfies \(\mathcal{F}_{\mathrm{G}}(\alpha) = 0\), we deduce from (31) that \(\alpha (\mathcal{F}_{\widehat{\mathrm{G}}} (f)) = 0\) for every \(f\in \mathrm{L}^{1}(\widehat{\mathrm{G}})\). As the image of \(\mathrm{L}^1 (\widehat{\mathrm{G}})\) by the Fourier transform is dense in \(\mathcal{C}_0(\mathrm{G})\) (corollary of Prop. 5 of II, p.~209), we therefore have \(\alpha = 0\) .

There exist functional spaces on G and \(\widehat{\mathbf{G}}\) other than \(\mathrm{L}^2 (\mathrm{G})\) and \(\mathrm{L}^2 (\widehat{\mathrm{G}})\) on which \(\mathcal{F}\) and \(\overline{\mathcal{F}}\) are inverse isomorphisms of each other. The following theorem gives an example. We denote by B(G) the vector subspace of \(\mathrm{L}^1 (\mathrm{G})\) consisting of the elements \(f\in \mathrm{L}^1 (\mathrm{G})\) such that \(\mathcal{F}_{\mathrm{G}}(f)\in \mathrm{L}^1 (\widehat{\mathrm{G}})\). It is a subalgebra of \(\mathrm{L}^1 (\mathrm{G})\). Indeed, let \(f\) and \(g\) be in B(G). We have \(f*g\in \mathrm{L}^1 (\mathrm{G})\) and \(\mathcal{F}_{\mathrm{G}}(f*g) = \mathcal{F}_{\mathrm{G}}(f)\mathcal{F}_{\mathrm{G}}(g)\in \mathrm{L}^1 (\widehat{\mathrm{G}})\) since \(\mathcal{F}_{\mathrm{G}}(f)\in \mathrm{L}^1 (\widehat{\mathrm{G}})\) and \(\mathcal{F}_{\mathrm{G}}(g)\in \mathcal{C}_0(\widehat{\mathrm{G}})\) .

THEOREM 3. --- The restriction of the Fourier transform to B(G) induces an isomorphism of vector spaces from B(G) onto B( \(\widehat{G}\) ), whose inverse is induced by the restriction to B( \(\widehat{G}\) ) of the Fourier cotransform.

Let \(f \in \mathrm{B}(\mathrm{G})\) . Denote by \(g = \mathcal{F}_{\mathrm{G}}(f)\) . We have \(g \in \mathrm{L}^1(\widehat{\mathrm{G}}) \cap \mathcal{C}_0(\widehat{\mathrm{G}}) \subset \mathrm{L}^1(\widehat{\mathrm{G}}) \cap \mathrm{L}^2(\widehat{\mathrm{G}})\) . Set \(f_1 = \overline{\mathcal{F}}_{\widehat{\mathrm{G}}} (g) \in \mathrm{L}^2(\mathrm{G})\). For every continuous function with compact support \(h \in \mathcal{K}(\widehat{\mathrm{G}})\) , one has \(h \in \mathrm{L}^1(\widehat{\mathrm{G}}) \cap \mathrm{L}^2(\widehat{\mathrm{G}})\) and

\[
\begin{array}{r l} & {\int_ {\mathrm{G}} f _ {1} (x) \mathcal {F} _ {\widehat {\mathrm{G}}} (h) (x) d x = \int_ {\mathrm{G}} \overline {{\mathcal {F}}} _ {\widehat {\mathrm{G}}} (g) (x) \overline {{\overline {{\mathcal {F}}} _ {\widehat {\mathrm{G}}} (\overline {{h}}) (x)}} d x} \\ & {\qquad = \int_ {\widehat {\mathrm{G}}} g (\chi) h (\chi) d \widehat {x} (\chi)} \\ & {\qquad = \int_ {\widehat {\mathrm{G}}} \mathcal {F} _ {\mathrm{G}} (f) (\chi) h (\chi) d \widehat {x} (\chi) = \int_ {\mathrm{G}} f (x) \mathcal {F} _ {\widehat {\mathrm{G}}} (h) (x) d x} \end{array}
\]

using the Plancherel theorem and formula (31). Since \(\mathcal{K}(\widehat{\mathrm{G}})\) is dense in \(\mathrm{L}^2 (\widehat{\mathrm{G}})\) , its image under the Fourier transform is dense in \(\mathrm{L}^2 (\mathrm{G})\) . Consequently, we have \(f_{1} = f\) in \(\mathrm{L}^{1}(\mathrm{G})\) ; this shows that \(g\in \mathrm{B}(\widehat{\mathrm{G}})\) .

The formula \(f_{1}=f\) means that the restriction to B(G) of the composition \(\overline{F}_{\widehat{G}}\circ F_{G}\) is the identity map of B(G). Interchanging the roles of G and \(\widehat{G}\) , we observe that \(F_{G}\circ\overline{F}_{\widehat{G}}\) is the identity map of B( \(\widehat{G}\) ), which completes the proof of the theorem.

COROLLARY 1. --- Let \(f \in \mathrm{L}^{1}(\mathrm{G})\) . Then \(f \in \mathrm{B}(\mathrm{G})\) if and only if \(f\) belongs to the image of the Fourier transform \(\mathcal{F}_{\widehat{\mathrm{G}}}\) on \(\mathrm{L}^{1}(\widehat{\mathrm{G}})\) . In particular, we have \(\mathrm{A}(\mathrm{G}) \subset \mathrm{B}(\mathrm{G})\) .

Theorem 3 proves that if \(f \in \mathrm{B}(\mathrm{G})\) , then \(f = \mathcal{F}_{\widehat{\mathrm{G}}}(\overline{\mathcal{F}}_{\mathrm{G}}(f))\) , where \(\overline{\mathcal{F}}_{\mathrm{G}}(f)\) belongs to \(\mathrm{L}^1(\widehat{\mathrm{G}})\). Conversely, if \(f = \mathcal{F}_{\widehat{\mathrm{G}}}(g)\) , where \(g \in \mathrm{L}^1(\widehat{\mathrm{G}})\) , then we have \(g \in \mathrm{B}(\widehat{\mathrm{G}})\) and therefore \(f \in \mathrm{B}(\mathrm{G})\) by the theorem. The last assertion then follows from prop. 11 of II, p.~217.

COROLLARY 2. --- The vector space B(G) is an algebra both for multiplication and for convolution. The Fourier transform interchanges convolution and multiplication in B(G) and B( \(\widehat{G}\) ).

We have already seen that B(G) is a subalgebra of L \(^{1}\) (G). On the other hand, if \(f\) and \(g\) belong to B(G), then \(fg \in \text{L}^{1}(\text{G})\) since \(f \in \text{L}^{1}(\text{G})\) and \(g\) belongs to the image of the Fourier transform on L \(^{1}\) ( \(\widehat{\text{G}}\) ) (corollary 1). Since there exist \(f_{1}\) and \(g_{1}\) in L \(^{1}\) ( \(\widehat{\text{G}}\) ) such that \(f = \mathcal{F}_{\widehat{\text{G}}} (f_{1})\) and \(g = \mathcal{F}_{\widehat{\text{G}}} (g_{1})\) (loc. cit.), we have \(fg = \mathcal{F}_{\widehat{\text{G}}} (f_{1} * g_{1})\) , and hence \(fg \in \text{B(G)}\) again by the preceding corollary.

PROPOSITION 14. — Let \(f\) and \(g\) be in \(\mathrm{L}^2 (\mathrm{G})\) . Then \(\mathcal{F}_{\mathrm{G}}(fg) = \mathcal{F}_{\mathrm{G}}(f)*\mathcal{F}_{\mathrm{G}}(g)\) .

The equality holds if \(f\) and \(g\) belong to B(G) (Corollary 2), and in particular if \(f\) and \(g\) belong to A(G) since A(G) ⊂ B(G) (Cor. 1). Since A(G) is dense in L²(G) (Cor. 1 of II, p.~212), it suffices to show that the two members of the equality are continuous functions of \((f,g)\in\mathrm{L}^{2}(\mathrm{G})\times\mathrm{L}^{2}(\mathrm{G})\) with values in \(\mathcal{C}_{0}(\widehat{\mathrm{G}})\) . Now the map \((f,g)\mapsto\mathcal{F}_{\mathrm{G}}(fg)\) is obtained by composing the continuous map \((f,g)\mapsto fg\) of \(\mathrm{L}^{2}(\mathrm{G})\times\mathrm{L}^{2}(\mathrm{G})\) into \(\mathrm{L}^{1}(\mathrm{G})\) and the Fourier transform \(\mathcal{F}_{\mathrm{G}}\) of \(\mathrm{L}^{1}(\mathrm{G})\) into \(\mathcal{C}_{0}(\widehat{\mathrm{G}})\) , which is also continuous. Likewise, the map \((f,g)\mapsto\mathcal{F}_{\mathrm{G}}(f)*\mathcal{F}_{\mathrm{G}}(g)\) is obtained by composing the continuous maps \((f,g)\mapsto(\mathcal{F}_{\mathrm{G}}(f),\mathcal{F}_{\mathrm{G}}(g))\) of \(\mathrm{L}^{2}(\mathrm{G})\times\mathrm{L}^{2}(\mathrm{G})\) into \(\mathrm{L}^{2}(\widehat{\mathrm{G}})\times\mathrm{L}^{2}(\widehat{\mathrm{G}})\) and \((h_{1},h_{2})\mapsto h_{1}*h_{2}\) of \(\mathrm{L}^{2}(\widehat{\mathrm{G}})\times\mathrm{L}^{2}(\widehat{\mathrm{G}})\) into \(\mathcal{C}_{0}(\widehat{\mathrm{G}})\) (INT, VIII, §4, no. 5, prop. 15).

Remark. --- See no. 9 and exercises 22 of II, p.~270 and 31 of II, p.~275 for other examples of function spaces on which the Fourier transform is an isomorphism, in the case of particular groups G.

